\documentclass[11pt]{article}
\usepackage{amsmath,amssymb,amsthm,mathtools}
\usepackage[margin=1in]{geometry}
\newtheorem{theorem}{Theorem}
\newtheorem{lemma}{Lemma}
\newtheorem{proposition}{Proposition}
\newtheorem{definition}{Definition}
\newtheorem{remark}{Remark}
\newtheorem{corollary}{Corollary}
\DeclareMathOperator{\Ran}{Ran}
\DeclareMathOperator{\Fix}{Fix}
\DeclareMathOperator{\tr}{tr}
\title{Step 84: Completed CND Gate for the Paired Weil Feature Package}
\author{Six Birds / Formed-Layer Membrane Project}
\date{}
\begin{document}
\maketitle

\section{Purpose}
The previous steps showed that trying to split the explicit formula into signed terms plus separated diagonal repairs is not the right carrier-level target.  A finite-rank pole/completion feature cannot pay a full diagonal on an infinite anti-invariant core, and a singular archimedean term is positive only as a paired difference form.  The right target is therefore a paired completed Weil feature representation.

This note formulates the gate for that target.  It does not prove RH.  It asks whether the completed zeta/Weil carrier can be represented, on the completed anti-invariant response core, by a positive paired feature package whose translation-invariant part is conditionally negative definite.

\section{Paired shift features}
Let \(\tau_a f(t)=f(t+a)\) on a log-side Hilbert space.  For a positive symmetric measure \(\mu\) satisfying the Levy integrability condition
\[
  \int_{\mathbb R}\min(a^2,1)\,d\mu(a)<\infty,
\]
define the paired shift form
\[
  q_\mu[f]=\int_{\mathbb R}\|\tau_a f-f\|_2^2\,d\mu(a).
\]
For suitable core functions,
\[
  q_\mu[f]=\int_{\mathbb R}\Psi_\mu(\xi)|\widehat f(\xi)|^2\,d\xi,
\]
where
\[
  \Psi_\mu(\xi)=2\int_{\mathbb R}(1-\cos(a\xi))\,d\mu(a).
\]

\begin{definition}[CND paired feature gate]
A translation-invariant symbol \(\Psi\) passes the conservative paired shift-feature gate if
\[
  \Psi(\xi)=c\xi^2+2\int_{\mathbb R}(1-\cos(a\xi))\,d\mu(a),
  \qquad c\ge0,
\]
for a positive symmetric Levy measure \(\mu\).  Equivalently, \(\Psi\) is a continuous negative definite even function with \(\Psi(0)=0\) and no killing term.
\end{definition}

\begin{remark}[Killing versus paired difference]
One may allow a killing term \(\kappa\|f\|^2\), corresponding to a symbol \(\kappa+\Psi(\xi)\).  Such a term is a genuine positive feature, but it is not a conservative shift-difference feature.  In the RH membrane route, any such killing/diagonal term must be supplied by an upstream-visible completion record, not introduced as free money after the fact.
\end{remark}

\section{Prime and gamma slots}
\subsection{Prime finite-cutoff slot}
For a cutoff \(L\), define positive weights
\[
  w_{n,L}=\frac{\Lambda(n)}{\sqrt n}\,\alpha_L(\log n)^2,
  \qquad \alpha_L\ge0.
\]
The prime finite-cutoff paired symbol is
\[
  \Psi_{{\rm pr},L}(\xi)
  =2\sum_{n\ge2}w_{n,L}\bigl(1-\cos(\xi\log n)\bigr).
\]
This passes the CND paired-feature gate at every finite cutoff, since it is generated by a finite positive atomic measure
\[
  \mu_{{\rm pr},L}=\sum_{n\ge2}w_{n,L}\delta_{\log n}.
\]
The unearned part is not positivity; it is matching this positive paired feature to the signed prime term in the completed explicit formula with correct tail accounting.

\subsection{Gamma slot}
For the zeta gamma factor, previous steps isolated the normalized archimedean paired symbol
\[
  \Psi_\infty(\xi)=\frac12\int_0^\infty
  \frac{e^{-t/4}}{1-e^{-t}}\left(1-\cos\frac{\xi t}{2}\right)dt.
\]
This is a paired shift feature with positive density
\[
  \kappa_\infty(t)=\frac12\frac{e^{-t/4}}{1-e^{-t}},\qquad t>0,
\]
and shift size \(t/2\).  Since \(\kappa_\infty(t)\sim (2t)^{-1}\) near zero and decays exponentially at infinity,
\[
  \int_0^\infty \min(t^2,1)\kappa_\infty(t)\,dt<\infty.
\]
Therefore the gamma slot passes the CND paired-feature gate as a paired difference form.

\begin{remark}[No separated gamma diagonal]
Although \(\Psi_\infty\) is a positive paired feature, the separated integral \(\int\kappa_\infty(t)dt\) diverges.  The gamma slot cannot be split into a free infinite diagonal and a negative correlation term.  Its positivity lives only in the paired expression.
\end{remark}

\section{Completed paired carrier gate}
The completed paired candidate has the schematic form
\[
  q_{\rm cmp}[f]
  =c\|f'\|_2^2+q_{{\rm pr},L}[f]+q_\infty[f]+q_{\rm pole}[f]+q_{\rm tail}[f]+e_{\rm br}[f].
\]
Here \(q_{\rm pole}\) is finite-rank or quotient/completion data, \(q_{\rm tail}\) is a positive tail/support feature or tail budget, and \(e_{\rm br}\) is a declared bridge defect if exact matching is not available.

\begin{definition}[Completed paired Weil gate]
The completed paired Weil gate passes on a core \(\mathcal C^-_{\log}\) if:
\begin{enumerate}
\item the translation-invariant part of \(q_{\rm cmp}\) is CND / Levy--Khinchine paired;
\item finite-rank pole/completion data is positive and null-mode legal;
\item tail/support terms are positive features or declared defects with fixed/exhaustive ledger control;
\item the zero-side form obeys
\[
  q_Z[f]\le q_{\rm cmp}[f]+e_{\rm br}[f],\qquad f\in\mathcal C^-_{\log};
\]
\item the core-to-full response-space gate is closable and exhaustive.
\end{enumerate}
\end{definition}

If the gate passes with \(e_{\rm br}=0\), then
\[
  V_Z^*V_Z\preceq W_{\rm cmp}^*W_{\rm cmp},
\]
and hence, by Douglas factorization,
\[
  V_Z=T W_{\rm cmp},\qquad \|T\|\le1.
\]

\section{No-go: positivity is weaker than paired-feature representability}
\begin{proposition}[Nonnegative does not imply paired shift feature]
A nonnegative even symbol need not pass the CND paired-feature gate.
\end{proposition}
\begin{proof}
The paired-feature gate requires conditional negative definiteness: for every finite set \(\xi_i\) and coefficients \(c_i\) with \(\sum_i c_i=0\),
\[
  \sum_{i,j}\overline{c_i}c_j\Psi(\xi_i-\xi_j)\le0.
\]
There are many nonnegative even functions that fail this condition, for example powers behaving like \(|\xi|^p\) with \(p>2\) near the origin.  Hence spectral positivity alone is not the same as a shift-difference feature representation.
\end{proof}

\section{RH implication}
The active RH construction target is now more precise than signed explicit-formula matching:
\[
  \boxed{\text{construct a completed paired CND/positive Weil carrier and prove }q_Z\le q_{\rm cmp}.}
\]
The prime and gamma slots have natural paired positive candidates.  What remains unearned is the exact completed matching:
\begin{itemize}
\item the signed prime term must match the paired prime feature without illegal separated diagonal accounting;
\item the signed gamma term must match the paired gamma feature with constants, pole terms, and null modes audited;
\item the pole/completion slot can repair finite-dimensional directions but not a full-core diagonal;
\item the tail/support slot must provide fixed or exhaustive ledger control;
\item the full core inequality must extend to the completed anti-invariant response space.
\end{itemize}

\section{Nonclaim boundary}
This note does not prove RH, does not prove the Weil--Douglas bridge, and does not prove that the actual completed zeta explicit formula has a paired CND feature representation.  It proves the gate and records that the prime and gamma candidate slots pass the paired-feature positivity test.  The remaining task is the signed-to-paired matching of the actual completed Weil carrier.

\end{document}
